\chapter{Gallbladder Scan (HIDA Scan)}

\section*{Overview}
A hepatobiliary iminodiacetic acid (HIDA) scan, or cholescintigraphy, uses a small radioactive tracer and a gamma camera to evaluate bile production and flow through the liver, gallbladder, and bile ducts.

\section*{Why It Is Done}
It helps evaluate suspected acute cholecystitis, bile-duct obstruction, bile leaks, impaired gallbladder emptying, or selected postoperative complications.

\section*{How It Is Performed}
The tracer is injected into a vein and sequential images are obtained as it passes through the hepatobiliary system. A medicine may be given to stimulate gallbladder contraction.

\section*{Preparation, Risks, and Results}
Fasting and medicine instructions must be followed. Mild injection discomfort and rare tracer reactions are possible; radiation exposure is low. Results assess tracer filling, drainage, and gallbladder ejection when measured.

\section*{References}
\begin{enumerate}
  \item MedlinePlus. Hepatobiliary Scan. U.S. National Library of Medicine; 2025.
  \item Current Medical Diagnosis and Treatment. McGraw-Hill; 2025.
\end{enumerate}
\clearpage
